% TikZ 图：感应电荷量与 I-t 图像面积
% Day 23：电磁感应
\begin{tikzpicture}[scale=0.92,>=Latex]
  \tikzset{
    axis/.style={->,thick},
    curve/.style={very thick,blue},
    shadearea/.style={fill=blue!18,draw=none},
    note/.style={font=\small,align=center}
  }

  \begin{scope}[xshift=-3.1cm]
    \draw[axis] (0,0) -- (2.8,0) node[right] {$t$};
    \draw[axis] (0,0) -- (0,2.1) node[above] {$I$};
    \path[shadearea] (0.35,0) -- (0.35,1.7)
      .. controls (0.65,1.95) and (1.05,1.95) .. (1.35,1.7)
      -- (1.35,0) -- cycle;
    \draw[curve] (0.35,0) -- (0.35,1.7)
      .. controls (0.65,1.95) and (1.05,1.95) .. (1.35,1.7)
      -- (1.35,0);
    \node[note] at (1.35,-0.48) {变化快\\电流大、时间短};
    \node[blue,font=\small] at (1.85,1.25) {$q=$面积};
  \end{scope}

  \begin{scope}[xshift=3.1cm]
    \draw[axis] (0,0) -- (2.8,0) node[right] {$t$};
    \draw[axis] (0,0) -- (0,2.1) node[above] {$I$};
    \path[shadearea] (0.35,0) -- (0.35,0.9)
      .. controls (0.95,1.10) and (1.95,1.10) .. (2.55,0.9)
      -- (2.55,0) -- cycle;
    \draw[curve] (0.35,0) -- (0.35,0.9)
      .. controls (0.95,1.10) and (1.95,1.10) .. (2.55,0.9)
      -- (2.55,0);
    \node[note] at (1.45,-0.48) {变化慢\\电流小、时间长};
    \node[blue,font=\small] at (1.85,0.63) {$q=$面积};
  \end{scope}

  \node[note] at (0,2.55) {若 $\Delta\Phi$ 相同且 $R$ 不变，两图面积相同};
  \node[note] at (0,-1.18) {$q=\bar I\Delta t=\dfrac{n|\Delta\Phi|}{R}$};
\end{tikzpicture}
